\section{Условие}
{\bfseries Цель работы:}
Приобретение практических навыков в управлении процессами в операционной системе и обеспечении обмена данными между процессами посредством каналов (pipe).

{\bfseries Задание:}
Составить и отладить программу на языке Си, осуществляющую работу с процессами и взаимодействие между ними в одной из двух операционных систем. В результате работы программа (основной процесс) должен создать для решения задачи один или несколько дочерних процессов. Взаимодействие между процессами осуществляется через системные сигналы/события и/или каналы (pipe). Необходимо обрабатывать системные ошибки, которые могут возникнуть в результате работы.

{\bfseries Вариант:} 16

{\bfseries Описание варианта 16:}
Родительский процесс создает дочерний процесс. Первой строкой пользователь в консоль родительского процесса вводит имя файла, которое будет использовано для открытия File с таким именем на запись. Родительский процесс принимает от пользователя строки произвольной длины и пересылает их в pipe1. Процесс child проверяет строки на валидность правилу. Если строка соответствует правилу, то она выводится в стандартный поток вывода дочернего процесса, иначе в pipe2 выводится информация об ошибке. Родительский процесс полученные от child ошибки выводит в стандартный поток вывода. Правило проверки: строка должна оканчиваться на «.» или «;».

